% --- Archon physics formalization source begin ---
% archon:physics
% archon:covers PhyXMiniProblems/problem_phyx_mini_0226.lean
% archon:source-report reports/phyx_mini/problem_phyx_mini_0226.source.json

\chapter{Physics problem phyx\_mini\_0226}
\label{ch:PhyXMiniProblems_problem_phyx_mini_0226}

\paragraph{Problem source.}
\#\# Physical scenario

A \textbackslash{}( 0.250\textbackslash{},\textbackslash{}mathrm\{kg\} \textbackslash{}) block resting on a frictionless, horizontal surface is attached to a spring whose force constant is \textbackslash{}( 83.8\textbackslash{},\textbackslash{}mathrm\{N/m\} \textbackslash{}) as in figure. A horizontal force \textbackslash{}( \textbackslash{}vec\{F\} \textbackslash{}) causes the spring to stretch a distance of \textbackslash{}( 5.46\textbackslash{},\textbackslash{}mathrm\{cm\} \textbackslash{}) from its equilibrium position.

\#\# Question

What is the total energy stored in the system when the spring is stretched?

\#\# Answer choices

- A: A: 0.325J

- B: B: 0.225J

- C: C: 0.425J

- D: D: 0.125J

\#\# Figure caption (auxiliary; use the image as primary evidence)

The image depicts a physics scenario involving a spring and a block. The scene consists of a vertical wall on the left side to which a spring is attached. The spring, which is shown in a compressed state, is connected to a blue block on the right. 

The block is situated on a horizontal surface, and to the right of the block, there is a blue arrow indicating a force. This arrow is labeled with the vector notation "F" with an arrow on top, suggesting the direction of the force is towards the right, away from the spring and wall. The elements demonstrate concepts such as force, compression of the spring, and potential motion of the block due to the force applied.

\#\# Dataset metadata

Work and Energy; Physical Model Grounding Reasoning

\paragraph{Recorded answer/context.}
D: D: 0.125J

\paragraph{Figure/image path.}
/root/proposal\_for\_physic/science-mango/phyx\_mini\_run/phyx\_data/test\_image/226.png

\paragraph{Formalization target.}
create a compiling Lean file with sorry bodies at `PhyXMiniProblems/problem\_phyx\_mini\_0226.lean`. The Lean declarations must preserve the physical quantities, dimensions or dimensional roles, figure labels, governing-law hypotheses, and final relation expressed by this problem.
Use Mathlib/Physlib names found through LeanExplore where available. If a physics API is missing, introduce faithful local abstractions rather than scalar placeholder aliases.

\begin{theorem}[Physics formalization target]
\label{thm:physics:phyx_mini_0226:target}
\lean{PhyXMiniProblems.ProblemPhyXMini0226.problem_phyx_mini_0226}
\uses{def:physics:phyx-mini-0226:phyxminiproblems-problemphyxmini0226-matchessuppliedfigure, def:physics:phyx-mini-0226:phyxminiproblems-problemphyxmini0226-springblocksetup, def:physics:phyx-mini-0226:phyxminiproblems-problemphyxmini0226-matchesproblemdata, def:physics:phyx-mini-0226:phyxminiproblems-problemphyxmini0226-hasphysicalparameters, def:physics:phyx-mini-0226:phyxminiproblems-problemphyxmini0226-satisfieselasticspringmodel, def:physics:phyx-mini-0226:phyxminiproblems-problemphyxmini0226-recordedanswerchoice, def:physics:phyx-mini-0226:phyxminiproblems-problemphyxmini0226-matchesdisplayedenergy, def:physics:phyx-mini-0226:phyxminiproblems-problemphyxmini0226-isuniqueclosestdisplayedchoice}
The assigned autoformalize agent should translate this physics problem into Lean declarations in the covered file, with theorem and lemma proof bodies written as `by sorry`.
\end{theorem}
\begin{proof}
This is an autoformalization task, not a proof task. Produce statements that can later be proved by the physics prover without weakening the physical model.
\end{proof}

% --- Archon declaration topology begin ---
\section{Declaration topology}
\begin{definition}[Mass Quantity]
  \label{def:physics:phyx-mini-0226:phyxminiproblems-problemphyxmini0226-massquantity}
  \lean{PhyXMiniProblems.ProblemPhyXMini0226.MassQuantity}
  Physical mass, with dimension M
\end{definition}
\begin{proof}
  This is the declared model definition of Mass Quantity.
\end{proof}
\begin{definition}[Length Quantity]
  \label{def:physics:phyx-mini-0226:phyxminiproblems-problemphyxmini0226-lengthquantity}
  \lean{PhyXMiniProblems.ProblemPhyXMini0226.LengthQuantity}
  Physical displacement or length, with dimension L
\end{definition}
\begin{proof}
  This is the declared model definition of Length Quantity.
\end{proof}
\begin{definition}[Velocity Quantity]
  \label{def:physics:phyx-mini-0226:phyxminiproblems-problemphyxmini0226-velocityquantity}
  \lean{PhyXMiniProblems.ProblemPhyXMini0226.VelocityQuantity}
  Physical velocity, with dimension L T⁻¹
\end{definition}
\begin{proof}
  This is the declared model definition of Velocity Quantity.
\end{proof}
\begin{definition}[Spring Constant Quantity]
  \label{def:physics:phyx-mini-0226:phyxminiproblems-problemphyxmini0226-springconstantquantity}
  \lean{PhyXMiniProblems.ProblemPhyXMini0226.SpringConstantQuantity}
  Physical spring stiffness, with dimension force per length, equivalently M T⁻²
\end{definition}
\begin{proof}
  This is the declared model definition of Spring Constant Quantity.
\end{proof}
\begin{definition}[Force Quantity]
  \label{def:physics:phyx-mini-0226:phyxminiproblems-problemphyxmini0226-forcequantity}
  \lean{PhyXMiniProblems.ProblemPhyXMini0226.ForceQuantity}
  Physical force, with dimension M L T⁻²
\end{definition}
\begin{proof}
  This is the declared model definition of Force Quantity.
\end{proof}
\begin{definition}[Energy Quantity]
  \label{def:physics:phyx-mini-0226:phyxminiproblems-problemphyxmini0226-energyquantity}
  \lean{PhyXMiniProblems.ProblemPhyXMini0226.EnergyQuantity}
  Physical energy, with dimension M L² T⁻²
\end{definition}
\begin{proof}
  This is the declared model definition of Energy Quantity.
\end{proof}
\begin{definition}[mass In Kilograms]
  \label{def:physics:phyx-mini-0226:phyxminiproblems-problemphyxmini0226-massinkilograms}
  \lean{PhyXMiniProblems.ProblemPhyXMini0226.massInKilograms}
  \uses{def:physics:phyx-mini-0226:phyxminiproblems-problemphyxmini0226-massquantity}
  SI kilogram readout of a physical mass
\end{definition}
\begin{proof}
  This is the declared model definition of mass In Kilograms.
\end{proof}
\begin{definition}[length In Meters]
  \label{def:physics:phyx-mini-0226:phyxminiproblems-problemphyxmini0226-lengthinmeters}
  \lean{PhyXMiniProblems.ProblemPhyXMini0226.lengthInMeters}
  \uses{def:physics:phyx-mini-0226:phyxminiproblems-problemphyxmini0226-lengthquantity}
  SI metre readout of a physical displacement
\end{definition}
\begin{proof}
  This is the declared model definition of length In Meters.
\end{proof}
\begin{definition}[velocity In Meters Per Second]
  \label{def:physics:phyx-mini-0226:phyxminiproblems-problemphyxmini0226-velocityinmeterspersecond}
  \lean{PhyXMiniProblems.ProblemPhyXMini0226.velocityInMetersPerSecond}
  \uses{def:physics:phyx-mini-0226:phyxminiproblems-problemphyxmini0226-velocityquantity}
  SI metre-per-second readout of a physical velocity
\end{definition}
\begin{proof}
  This is the declared model definition of velocity In Meters Per Second.
\end{proof}
\begin{definition}[spring Constant In Newtons Per Meter]
  \label{def:physics:phyx-mini-0226:phyxminiproblems-problemphyxmini0226-springconstantinnewtonspermeter}
  \lean{PhyXMiniProblems.ProblemPhyXMini0226.springConstantInNewtonsPerMeter}
  \uses{def:physics:phyx-mini-0226:phyxminiproblems-problemphyxmini0226-springconstantquantity}
  SI newton-per-metre readout of a physical spring stiffness
\end{definition}
\begin{proof}
  This is the declared model definition of spring Constant In Newtons Per Meter.
\end{proof}
\begin{definition}[force In Newtons]
  \label{def:physics:phyx-mini-0226:phyxminiproblems-problemphyxmini0226-forceinnewtons}
  \lean{PhyXMiniProblems.ProblemPhyXMini0226.forceInNewtons}
  \uses{def:physics:phyx-mini-0226:phyxminiproblems-problemphyxmini0226-forcequantity}
  SI newton readout of a physical force magnitude
\end{definition}
\begin{proof}
  This is the declared model definition of force In Newtons.
\end{proof}
\begin{definition}[energy In Joules]
  \label{def:physics:phyx-mini-0226:phyxminiproblems-problemphyxmini0226-energyinjoules}
  \lean{PhyXMiniProblems.ProblemPhyXMini0226.energyInJoules}
  \uses{def:physics:phyx-mini-0226:phyxminiproblems-problemphyxmini0226-energyquantity}
  SI joule readout of a physical energy
\end{definition}
\begin{proof}
  This is the declared model definition of energy In Joules.
\end{proof}
\begin{definition}[Figure Side]
  \label{def:physics:phyx-mini-0226:phyxminiproblems-problemphyxmini0226-figureside}
  \lean{PhyXMiniProblems.ProblemPhyXMini0226.FigureSide}
  Left and right sides of the supplied spring-block diagram
\end{definition}
\begin{proof}
  This is the declared model definition of Figure Side.
\end{proof}
\begin{definition}[Horizontal Direction]
  \label{def:physics:phyx-mini-0226:phyxminiproblems-problemphyxmini0226-horizontaldirection}
  \lean{PhyXMiniProblems.ProblemPhyXMini0226.HorizontalDirection}
  Horizontal direction of the labeled force arrow
\end{definition}
\begin{proof}
  This is the declared model definition of Horizontal Direction.
\end{proof}
\begin{definition}[Surface Orientation]
  \label{def:physics:phyx-mini-0226:phyxminiproblems-problemphyxmini0226-surfaceorientation}
  \lean{PhyXMiniProblems.ProblemPhyXMini0226.SurfaceOrientation}
  Orientation of the support surface drawn beneath the block
\end{definition}
\begin{proof}
  This is the declared model definition of Surface Orientation.
\end{proof}
\begin{definition}[Spring Endpoint]
  \label{def:physics:phyx-mini-0226:phyxminiproblems-problemphyxmini0226-springendpoint}
  \lean{PhyXMiniProblems.ProblemPhyXMini0226.SpringEndpoint}
  The two physical objects joined by the spring in the figure
\end{definition}
\begin{proof}
  This is the declared model definition of Spring Endpoint.
\end{proof}
\begin{definition}[Spring Block Figure]
  \label{def:physics:phyx-mini-0226:phyxminiproblems-problemphyxmini0226-springblockfigure}
  \lean{PhyXMiniProblems.ProblemPhyXMini0226.SpringBlockFigure}
  \uses{def:physics:phyx-mini-0226:phyxminiproblems-problemphyxmini0226-figureside, def:physics:phyx-mini-0226:phyxminiproblems-problemphyxmini0226-horizontaldirection, def:physics:phyx-mini-0226:phyxminiproblems-problemphyxmini0226-surfaceorientation, def:physics:phyx-mini-0226:phyxminiproblems-problemphyxmini0226-springendpoint}
  Raw labels and geometry visible in the supplied image
\end{definition}
\begin{proof}
  This is the declared model definition of Spring Block Figure.
\end{proof}
\begin{definition}[Matches Supplied Figure]
  \label{def:physics:phyx-mini-0226:phyxminiproblems-problemphyxmini0226-matchessuppliedfigure}
  \lean{PhyXMiniProblems.ProblemPhyXMini0226.MatchesSuppliedFigure}
  \uses{def:physics:phyx-mini-0226:phyxminiproblems-problemphyxmini0226-springblockfigure}
  Primary-image evidence: wall and spring to the left, block to the right, a horizontal support, and the arrow F directed rightward away from the wall
\end{definition}
\begin{proof}
  This is the declared model definition of Matches Supplied Figure.
\end{proof}
\begin{definition}[Spring Block Setup]
  \label{def:physics:phyx-mini-0226:phyxminiproblems-problemphyxmini0226-springblocksetup}
  \lean{PhyXMiniProblems.ProblemPhyXMini0226.SpringBlockSetup}
  \uses{def:physics:phyx-mini-0226:phyxminiproblems-problemphyxmini0226-massquantity, def:physics:phyx-mini-0226:phyxminiproblems-problemphyxmini0226-lengthquantity, def:physics:phyx-mini-0226:phyxminiproblems-problemphyxmini0226-velocityquantity, def:physics:phyx-mini-0226:phyxminiproblems-problemphyxmini0226-springconstantquantity, def:physics:phyx-mini-0226:phyxminiproblems-problemphyxmini0226-forcequantity, def:physics:phyx-mini-0226:phyxminiproblems-problemphyxmini0226-energyquantity, def:physics:phyx-mini-0226:phyxminiproblems-problemphyxmini0226-springblockfigure}
  The physical spring-block system and the quantities named by the problem. extensionVectorInMeters is the one-dimensional SI-coordinate readout used by Physlib. The physical extension and stored energy remain dimensionful fields, and no numerical answer or answer-choice label is stored in this structure
\end{definition}
\begin{proof}
  This is the declared model definition of Spring Block Setup.
\end{proof}
\begin{definition}[Matches Problem Data]
  \label{def:physics:phyx-mini-0226:phyxminiproblems-problemphyxmini0226-matchesproblemdata}
  \lean{PhyXMiniProblems.ProblemPhyXMini0226.MatchesProblemData}
  \uses{def:physics:phyx-mini-0226:phyxminiproblems-problemphyxmini0226-massinkilograms, def:physics:phyx-mini-0226:phyxminiproblems-problemphyxmini0226-lengthinmeters, def:physics:phyx-mini-0226:phyxminiproblems-problemphyxmini0226-velocityinmeterspersecond, def:physics:phyx-mini-0226:phyxminiproblems-problemphyxmini0226-springconstantinnewtonspermeter, def:physics:phyx-mini-0226:phyxminiproblems-problemphyxmini0226-forceinnewtons, def:physics:phyx-mini-0226:phyxminiproblems-problemphyxmini0226-springblocksetup}
  Numerical measurements and qualitative conditions stated in the problem. The extension is converted from 5.46 cm to metres. Initial rest and zero friction are properties of the setup, not assumptions about the requested stored-energy value
\end{definition}
\begin{proof}
  This is the declared model definition of Matches Problem Data.
\end{proof}
\begin{definition}[Has Physical Parameters]
  \label{def:physics:phyx-mini-0226:phyxminiproblems-problemphyxmini0226-hasphysicalparameters}
  \lean{PhyXMiniProblems.ProblemPhyXMini0226.HasPhysicalParameters}
  \uses{def:physics:phyx-mini-0226:phyxminiproblems-problemphyxmini0226-massinkilograms, def:physics:phyx-mini-0226:phyxminiproblems-problemphyxmini0226-lengthinmeters, def:physics:phyx-mini-0226:phyxminiproblems-problemphyxmini0226-springconstantinnewtonspermeter, def:physics:phyx-mini-0226:phyxminiproblems-problemphyxmini0226-forceinnewtons, def:physics:phyx-mini-0226:phyxminiproblems-problemphyxmini0226-springblocksetup}
  Positivity conditions for the physical quantities supplied by the setup
\end{definition}
\begin{proof}
  This is the declared model definition of Has Physical Parameters.
\end{proof}
\begin{definition}[Satisfies Elastic Spring Model]
  \label{def:physics:phyx-mini-0226:phyxminiproblems-problemphyxmini0226-satisfieselasticspringmodel}
  \lean{PhyXMiniProblems.ProblemPhyXMini0226.SatisfiesElasticSpringModel}
  \uses{def:physics:phyx-mini-0226:phyxminiproblems-problemphyxmini0226-massinkilograms, def:physics:phyx-mini-0226:phyxminiproblems-problemphyxmini0226-lengthinmeters, def:physics:phyx-mini-0226:phyxminiproblems-problemphyxmini0226-springconstantinnewtonspermeter, def:physics:phyx-mini-0226:phyxminiproblems-problemphyxmini0226-energyinjoules, def:physics:phyx-mini-0226:phyxminiproblems-problemphyxmini0226-springblocksetup}
  The one-dimensional Hooke-law energy model in coherent SI readouts. Physlib's oscillator packages the positive mass and spring constant, while potentialEnergy is its general law U = (1/2) k x². The last field says that the physical energy stored by stretching the spring has that SI readout. No problem-specific energy value or answer choice occurs in this interface
\end{definition}
\begin{proof}
  This is the declared model definition of Satisfies Elastic Spring Model.
\end{proof}
\begin{definition}[Answer Choice]
  \label{def:physics:phyx-mini-0226:phyxminiproblems-problemphyxmini0226-answerchoice}
  \lean{PhyXMiniProblems.ProblemPhyXMini0226.AnswerChoice}
  Labels printed beside the four candidate energies
\end{definition}
\begin{proof}
  This is the declared model definition of Answer Choice.
\end{proof}
\begin{definition}[energy In Joules]
  \label{def:physics:phyx-mini-0226:phyxminiproblems-problemphyxmini0226-answerchoice-energyinjoules}
  \lean{PhyXMiniProblems.ProblemPhyXMini0226.AnswerChoice.energyInJoules}
  \uses{def:physics:phyx-mini-0226:phyxminiproblems-problemphyxmini0226-energyinjoules, def:physics:phyx-mini-0226:phyxminiproblems-problemphyxmini0226-answerchoice}
  Energy in joules printed beside each answer label
\end{definition}
\begin{proof}
  This is the declared model definition of energy In Joules.
\end{proof}
\begin{definition}[recorded Answer Choice]
  \label{def:physics:phyx-mini-0226:phyxminiproblems-problemphyxmini0226-recordedanswerchoice}
  \lean{PhyXMiniProblems.ProblemPhyXMini0226.recordedAnswerChoice}
  \uses{def:physics:phyx-mini-0226:phyxminiproblems-problemphyxmini0226-answerchoice}
  The answer label recorded in the supplied dataset metadata
\end{definition}
\begin{proof}
  This is the declared model definition of recorded Answer Choice.
\end{proof}
\begin{definition}[Matches Displayed Energy]
  \label{def:physics:phyx-mini-0226:phyxminiproblems-problemphyxmini0226-matchesdisplayedenergy}
  \lean{PhyXMiniProblems.ProblemPhyXMini0226.MatchesDisplayedEnergy}
  \uses{def:physics:phyx-mini-0226:phyxminiproblems-problemphyxmini0226-energyinjoules, def:physics:phyx-mini-0226:phyxminiproblems-problemphyxmini0226-springblocksetup, def:physics:phyx-mini-0226:phyxminiproblems-problemphyxmini0226-answerchoice}
  Agreement with a displayed energy to the nearest millijoule. The half- millijoule tolerance avoids identifying the exact value from the supplied three-significant-figure inputs with the rounded display 0.125 J
\end{definition}
\begin{proof}
  This is the declared model definition of Matches Displayed Energy.
\end{proof}
\begin{definition}[Is Unique Closest Displayed Choice]
  \label{def:physics:phyx-mini-0226:phyxminiproblems-problemphyxmini0226-isuniqueclosestdisplayedchoice}
  \lean{PhyXMiniProblems.ProblemPhyXMini0226.IsUniqueClosestDisplayedChoice}
  \uses{def:physics:phyx-mini-0226:phyxminiproblems-problemphyxmini0226-energyinjoules, def:physics:phyx-mini-0226:phyxminiproblems-problemphyxmini0226-springblocksetup, def:physics:phyx-mini-0226:phyxminiproblems-problemphyxmini0226-answerchoice}
  The selected energy is strictly closer than every other displayed choice
\end{definition}
\begin{proof}
  This is the declared model definition of Is Unique Closest Displayed Choice.
\end{proof}
% --- Archon declaration topology end ---
% --- Archon physics formalization source end ---
